\subsection{LC High Pass Filter}

\subsubsection{Overview}
Constructed from an inductor and capacitor in series, a LC
High Pass Filter allows high frequency signals to pass while
attenuating low frequency signals. It is a second order
filter and can provide a steeper roll-off (-40 dB/decade) than
single pole filters (-20 dB/decade).

\subsubsection{Schematic}
\begin{center}
  \begin{circuitikz}
    % Paths, nodes and wires:
	  \draw (5.5, 5.25) to[american voltage source, l_={$V_{ee}$}] (5.5, 6.5);
	  \node[ground] at (5.5, 8){};
	  \node[ground, xscale=-1, yscale=-1] at (5.5, 6.5){};
	  \draw (7, 7.75) -- (7, 6.75);
	  \draw (7, 5.25) -| (7, 5);
	  \draw (5.5, 5) -- (7, 5);
	  \draw (5.5, 5) -| (5.5, 5.25);
	  \draw (5.5, 9.25) -| (5.5, 9.5);
	  \draw (5.5, 9.5) -- (7, 9.5);
	  \draw (7, 9.5) -| (7, 9.25);
	  \draw (7, 7.25) -- (7.75, 7.25);
	  \draw (5.5, 8) to[sinusoidal voltage source, l_={$V_{in}$}] (5.5, 9.25);
	  \node[shape=rectangle, minimum width=0.965cm, minimum height=0.465cm](N1) at (8.25, 7.25){} node[anchor=center] at (N1.text){$V_{out}$};
	  \node[ground] at (9.5, 8){};
	  \node[ground, xscale=-1, yscale=-1] at (9.5, 5.25){};
	  \draw (11, 7.75) -- (11, 6.75);
	  \draw (11, 5.25) -| (11, 5);
	  \draw (9.5, 5) -- (11, 5);
	  \draw (9.5, 5) -| (9.5, 5.25);
	  \draw (9.5, 9.25) -| (9.5, 9.5);
	  \draw (9.5, 9.5) -- (11, 9.5);
	  \draw (11, 9.5) -| (11, 9.25);
	  \draw (11, 7.25) -- (11.75, 7.25);
	  \draw (9.5, 8) to[sinusoidal voltage source, l_={$V_{in}$}] (9.5, 9.25);
	  \node[shape=rectangle, minimum width=0.965cm, minimum height=0.465cm](N2) at (12.25, 7.25){} node[anchor=center] at (N2.text){$V_{out}$};
	  \node[shape=rectangle, minimum width=3.215cm, minimum height=0.465cm] at (6.125, 4.25){} node[anchor=north, align=center, text width=2.827cm, inner sep=6pt] at (6.125, 4.5){Full Schematic};
	  \node[shape=rectangle, minimum width=2.715cm, minimum height=0.465cm] at (10.375, 4.25){} node[anchor=north, align=center, text width=2.327cm, inner sep=6pt] at (10.375, 4.5){Small-Signal\\Schematic};
	  \draw (7, 6.75) to[american inductor, l={$L_1$}] (7, 5.25);
	  \draw (11, 6.75) to[american inductor, l={$L_1$}] (11, 5.25);
	  \draw (7, 7.75) to[capacitor, l_={$C_1$}] (7, 9.25);
	  \draw (11, 7.75) to[capacitor, l_={$C_1$}] (11, 9.25);
  \end{circuitikz}
\end{center}

\subsubsection{Transfer Function}
Applying KCL at $V_{out}$ on Full Schematic, we get:
\[
  (V_{out} - V_{in})sC_1 + \frac{(V_{out} - V_{ee})}{sL_1} = 0
\]
\paragraph{Superposition:}
\begin{enumerate}
  \item Due to \(V_{in}\) alone (\(V_{ee} = 0V\)):
    \begin{align*}
      (V_{out} - V_{in})sC_1 + \frac{(V_{out} - 0)}{sL_1} &= 0 \\
      V_{out}sC_1 - V_{in}sC_1 + \frac{V_{out}}{sL_1} &= 0 \\
      V_{out}(sC_1 + \frac{1}{sL_1}) - V_{in}sC_1 &= 0 \\
      V_{out}(sC_1 + \frac{1}{sL_1}) &= V_{in}sC_1 \\
      V_{out} &= \frac{V_{in}sC_1}{sC_1 + \frac{1}{sL_1}} \\
      V_{out} &= \frac{V_{in}sC_1}{\frac{s^2C_1L_1 + 1}{sL_1}} \\
      V_{out} &= \frac{V_{in}s^2C_1L_1}{s^2C_1L_1 + 1} \\
      \boxed{H_1(s) = \frac{N_1(s)}{D_1(s)} = \frac{V_{out_1}}{V_{in}} = \frac{s^2C_1L_1}{s^2C_1L_1 + 1}}
    \end{align*}
  \item Due to \(V_{ee}\) alone (\(V_{in} = 0V\)):
    \begin{align*}
      (V_{out} - 0)sC_1 + \frac{(V_{out} - V_{ee})}{sL_1} &= 0 \\
      V_{out}sC_1 + \frac{V_{out}}{sL_1} - \frac{V_{ee}}{sL_1} &= 0 \\
      V_{out}(sC_1 + \frac{1}{sL_1}) - \frac{V_{ee}}{sL_1} &= 0 \\
      V_{out}(sC_1 + \frac{1}{sL_1}) &= \frac{V_{ee}}{sL_1} \\
      V_{out} &= \frac{\frac{V_{ee}}{sL_1}}{(sC_1 + \frac{1}{sL_1})} \\
      V_{out} &= \frac{\frac{V_{ee}}{sL_1}}{\frac{s^2C_1L_1 + 1}{sL_1}} \\
      V_{out} &= \frac{V_{ee}}{s^2C_1L_1 + 1} \\
      \boxed{H_2(s) = \frac{N_2(s)}{D_2(s)} = \frac{V_{out_2}}{V_{ee}} = \frac{1}{s^2C_1L_1 + 1}}
    \end{align*}
\end{enumerate}
\paragraph{Total Output:}
\begin{align*}
  V_{out} &= V_{out_1} + V_{out_2} \\
  V_{out} &= H_1(s)V_{in} + H_2(s)V_{ee} \\
  V_{out} &= \frac{s^2C_1L_1}{s^2C_1L_1 + 1}V_{in} + \frac{1}{s^2C_1L_1 + 1}V_{ee} \\
\end{align*}

\subsubsection{Analysis}
\begin{enumerate}
  \item \textbf{Poles (\(D(s) = 0\)):} \\
    For both \(H_1(s)\) and \(H_2(s)\):
    \begin{align*}
      s^2C_1L_1 + 1 &= 0 \\
      s^2C_1L_1 &= -1 \\
      s^2 &= -\frac{1}{C_1L_1} \\
    \end{align*}
    Complex conjugate poles on the imaginary axis indicate
    a marginally stable second-order system. Frequency is given
    by:
    \begin{align*}
      s &= -2 \pi f_p \\
      (-2 \pi f_p)^2 &= -\frac{1}{C_1L_1} \\
      4 \pi^2 f_p^2 &= \frac{1}{C_1L_1} \\
      f_p &= \frac{1}{2 \pi \sqrt{C_1L_1}}
    \end{align*}
  \item \textbf{Zeroes (\(N(s) = 0\)):} \\
    For \(H_1(s)\):
    \begin{align*}
        N_1(s) = s^2C_1L_1 = 0 \\
        s_z &= 0 \quad \text{[let \(s = s_z\)]} \\
        f_z &= \frac{|s_z|}{2\pi} = 0 Hz
    \end{align*}
    Zero at the origin for \(H_1(s)\). \\
    For \(H_2(s)\):
    \begin{align*}
        N_2(s) = 1 \neq 0
    \end{align*}
    No zeroes for \(H_2(s)\).
  \item \textbf{DC Gain \(A_{DC}\):} \\
    For \(H_1(s)\):
    \begin{align*}
      H_1(0) = \frac{0^2C_1L_1}{0^2C_1L_1 + 1} = 0
    \end{align*}
    Zero DC gain for \(H_1(s)\). \\
    For \(H_2(s)\):
    \begin{align*}
      H_2(0) = \frac{1}{0^2C_1L_1 + 1} = 1
    \end{align*}
    Unity DC gain for \(H_2(s)\).
  \item \textbf{Gain Bandwidth Product (GBW):}
    We use the DC gain from \(H_1(s)\) since it is the
    high pass characteristic:
    \begin{align*}
      GBW &= A_{DC} \cdot f_p \\
      &= 0 \cdot \frac{1}{2 \pi \sqrt{C_1L_1}} \\
      \boxed{GBW = 0}
    \end{align*}
  \item \textbf{Impedances:}
    Same as the LC Low Pass Filter.
\end{enumerate}

\subsubsection{Question}
Design a LC High Pass Filter with a cutoff frequency of
\(f_c = 10 kHz\). Choose standard component values for \(L_1\)
and \(C_1\).
\paragraph{Solution:}
\begin{enumerate}
  \item Choosing standard value \(C_1 = 10 nF\).
  \item Calculating $L_1$:
    \begin{align*}
      f_c &= \frac{1}{2 \pi \sqrt{C_1L_1}} \\
      L_1 &= \frac{1}{(2 \pi f_c)^2 C_1} \\
      L_1 &= \frac{1}{(2 \pi \cdot 10 \times 10^3)^2 \cdot 10 \times 10^{-9}} \\
      L_1 &\approx 25.33 mH
    \end{align*}
    Choose standard value \(L_1 = 27 mH\).
  \item Verifying $f_p$ with chosen values:
    \begin{align*}
      f_p &= \frac{1}{2 \pi \sqrt{C_1L_1}} \\
      f_p &= \frac{1}{2 \pi \sqrt{0.01 \times 10^{-6} \cdot 27 \times 10^{-3}}} \\
      f_p &\approx 9.685 kHz
    \end{align*}
    Close to desired pole frequency.
\end{enumerate}
